% =====================================================================
%  Section: Side-to-side tower-mode excitation in the high-fidelity
%           OpenFAST wind-turbine model
%  Required in preamble:
%    \usepackage{amsmath, amssymb}
%    \usepackage{graphicx}
%    \usepackage{booktabs}
%    \usepackage{siunitx}
%  Adjust the figure path (fig/...) to your own.
% =====================================================================

\section{Side-to-Side Tower-Mode Excitation in the OpenFAST Model}
\label{sec:openfast-excitation}

Before adopting the reduced two-mass model of
\S\ref{sec:frequency-analysis}, a series of experiments was carried out on the
high-fidelity aeroelastic wind turbine, running OpenFAST as a Functional
Mock-up Unit (FMU) co-simulated with the electrical grid. The objective was to
excite the lightly damped first side-to-side (SS) tower mode from the
electrical side, i.e.\ to demonstrate an electro-mechanical interaction in the
detailed model. This section documents what was attempted, why the electrical
excitation path turned out to be blocked by the turbine controller, and how the
mode was ultimately excited through an aerodynamic path instead. The negative
result directly motivates the pivot to the simplified model.

\subsection{Model and co-simulation setup}
\label{subsec:of-setup}

The turbine is the IEA \SI{15}{\mega\watt} reference machine, exported from
OpenFAST as an FMU and coupled to the power-system model at a fixed
co-simulation step of $\Delta t = \SI{0.01}{\second}$. The relevant
configuration was:
\begin{itemize}
  \item \textbf{ElastoDyn}: only the structural degrees of freedom of interest
  enabled --- the first side-to-side tower DOF (\texttt{TwSSDOF1}), the
  generator DOF (\texttt{GenDOF}), and the rotor set spinning at
  \SI{5.70}{rpm}; blade pitch fixed and all other DOFs disabled to isolate the
  SS response.
  \item \textbf{ServoDyn}: the generator-torque and pitch control are handled by
  the ROSCO controller through the Bladed-style DLL interface
  ($\texttt{VSContrl}=5$, $\texttt{PCMode}=5$).
  \item \textbf{Hydrodynamics/substructure disabled} ($\texttt{CompHydro}=0$,
  $\texttt{CompSub}=0$): no wave loading, so the SS response is driven purely by
  the excitation under study.
  \item The FMU exposes the nacelle side-to-side acceleration
  (\texttt{YawBrTAyp}) as the SS observable, together with the generator torque
  (\texttt{GenTq}) and blade pitch (\texttt{BldPitch}).
\end{itemize}

\paragraph{The side-to-side tower mode.}
A free-decay (ring-down) test identified the first SS tower mode at
\begin{equation}
  f_{\mathrm{SS}} \approx \SI{0.234}{\hertz}
  \quad(\text{precisely } \SI{0.23411}{\hertz}),
  \qquad \zeta_{\mathrm{SS}} \approx \SI{0.5}{\percent},
  \qquad \tau_{\mathrm{SS}} \approx \SI{136}{\second},
  \label{eq:ss-mode}
\end{equation}
i.e.\ a very lightly damped, essentially aerodynamically decoupled lateral mode.
Its low damping is what makes a sustained resonance possible in principle; the
long ring-down time constant $\tau_{\mathrm{SS}}$ also means transients persist
for many minutes, a fact that complicates the interpretation of the results
(see below).

\subsection{Attempt 1: excitation through the electrical/torque command}
\label{subsec:of-electrical}

The first and most direct idea was to modulate the electrical loading of the
generator at the SS frequency and let the resulting torque ripple drive the
tower laterally. Two command channels exposed by the FMU were tried:
\begin{enumerate}
  \item a commanded electrical power \texttt{ElecPwrCom}, modulated by
  $\pm\SI{20}{\percent}$ at \SI{0.236}{\hertz};
  \item a commanded generator torque \texttt{GenSpdOrTrq}.
\end{enumerate}
The same tests were repeated below rated (Region~2, \SI{8}{\metre\per\second})
and above rated (Region~3), where in Region~2 the pitch remained at
\SI{0}{\degree} with zero ripple, confirming the operating region.

\paragraph{Result: the commands are rejected.}
A baseline run \emph{without} any modulation was compared against the modulated
runs. The generator-torque ripple was found to be \emph{bit-identical} in both
cases (e.g.\ $\pm\SI{144.96}{\kilo\newton\metre}$ in Region~2, with and without
\texttt{ElecPwrCom} modulation), and the SS acceleration envelope was
indistinguishable. In other words, the external power and torque commands have
\emph{zero} effect on the simulated torque. Table~\ref{tab:of-electrical}
summarises the comparison.

\begin{table}[t]
  \centering
  \caption{Electrical/torque command test (Region~2). The modulated and
  baseline cases are identical, showing the commands are ignored.}
  \label{tab:of-electrical}
  \begin{tabular}{lcc}
    \toprule
     & \texttt{ElecPwrCom} $\pm\SI{20}{\percent}$ & Baseline (no mod.) \\
    \midrule
    Mean \texttt{GenTq} [\si{\kilo\newton\metre}]   & 11\,179 & 11\,179 \\
    \texttt{GenTq} ripple [\si{\kilo\newton\metre}] & $\pm 144.96$ & $\pm 144.96$ \\
    Blade-pitch ripple                              & 0 & 0 \\
    SS envelope (decaying)                          & $0.204 \to 0.092$ & $0.204 \to 0.092$ \\
    \bottomrule
  \end{tabular}
\end{table}

\paragraph{Why: ROSCO owns the torque.}
The cause is the control architecture. With $\texttt{VSContrl}=5$ the generator
torque is computed \emph{internally} by ROSCO from the measured generator speed,
in Region~2 following the standard optimal $K\omega^2$ law,
\begin{equation}
  T_{\mathrm{gen}} = K_{\mathrm{opt}}\,\omega_{g}^{2},
  \label{eq:komega2}
\end{equation}
and does not read the external \texttt{ElecPwrCom}/\texttt{GenSpdOrTrq} ports.
The $\pm\SI{145}{\kilo\newton\metre}$ ripple that is observed is therefore not a
response to the command at all, but ROSCO's own speed-reaction to the start-up
transient. The only remaining physical grid$\to$SS path is indirect (a grid
event perturbs the electrical torque, which changes the speed, to which ROSCO
then reacts), but this coupling was found to be negligible: the SS response to
grid disturbances was roughly \SI{0.012}{\metre\per\second\squared} steady
against a \SI{0.225}{\metre\per\second\squared} start-up transient. The
conclusion is unambiguous: \emph{the SS mode cannot be excited through the
electrical/torque command channel in this FMU.}

\subsection{Attempt 2: bypassing ROSCO via \texttt{VSContrl}=4}
\label{subsec:of-vscontrl4}

To open the external-torque path, ServoDyn was reconfigured to
$\texttt{VSContrl}=4$ (external/Simulink torque), which would let the
grid-power-derived torque become the actual generator torque. The OpenFAST
input files are read at FMU run time, so this change requires no FMU rebuild.
However, the FMU \emph{crashes} under $\texttt{VSContrl}=4$ (an access-violation
error in the FMI call), so this route was abandoned and the controller reverted
to $\texttt{VSContrl}=5$. The electrical path thus remained closed.

\subsection{Attempt 3: aerodynamic excitation via oscillating wind direction}
\label{subsec:of-winddir}

Since the torque channel is governed by ROSCO, the excitation was instead
introduced \emph{upstream of the controller}, through the aerodynamics. A
uniform (hub-height) wind input file was generated in which the wind
\emph{direction} column oscillates sinusoidally at the SS frequency,
\begin{equation}
  \phi_{\mathrm{wind}}(t) = \hat{\phi}\,\sin\!\big(2\pi f_{\mathrm{SS}}\, t\big),
  \label{eq:winddir}
\end{equation}
while the wind speed is held constant. The oscillating direction produces a
periodic yaw misalignment, hence a lateral aerodynamic force on the rotor that
drives the tower side-to-side \emph{directly}, entirely bypassing the ROSCO
torque channel.

\paragraph{First run (amplitude \SI{4}{\degree}, \SI{0.234}{\hertz}).}
In contrast to every torque/power run --- where the SS envelope decayed
monotonically towards the ring-down floor --- the SS response here did
\emph{not} decay. After the initial start-up transient died away (a dip around
$t\approx\SI{110}{\second}$), the amplitude \emph{rebuilt} to a sustained
$\approx\SI{0.21}{\metre\per\second\squared}$. The re-growth after the dip is
the signature of sustained resonant forcing.

\paragraph{Delayed-start run (amplitude \SI{10}{\degree}, onset at
\SI{150}{\second}).}
Because the first run started the sinusoid at $t=0$, simultaneously with the
start-up transient (both near \SI{0.2}{\metre\per\second\squared} at the same
frequency), the two could not be visually separated. The experiment was
therefore repeated with the directional forcing delayed to
$t=\SI{150}{\second}$. This produced the clean, textbook picture: the SS
acceleration first decays to $\approx\SI{0.08}{\metre\per\second\squared}$ by
$t=\SI{150}{\second}$, and then \emph{grows monotonically} to
$\pm\SI{0.42}{\metre\per\second\squared}$ by $t=\SI{400}{\second}$, still
rising --- a classic resonance build-up (Figure~\ref{fig:ss-buildup}). This
confirms that the wind-direction oscillation at $f_{\mathrm{SS}}$ excites the
side-to-side mode.

\begin{figure}[t]
  \centering
  % \includegraphics[width=0.8\linewidth]{fig/ss_winddir_buildup.png}
  \caption{Side-to-side nacelle acceleration under wind-direction forcing at
  $f_{\mathrm{SS}}=\SI{0.234}{\hertz}$ with forcing onset at
  $t=\SI{150}{\second}$. The response decays to the ring-down floor, then builds
  up resonantly once the directional forcing begins.}
  \label{fig:ss-buildup}
\end{figure}

\subsection{Effect of the second side-to-side DOF}
\label{subsec:of-dof2}

Enabling the second SS tower DOF (\texttt{TwSSDOF2}) and repeating the forcing
markedly reduced the response, from $\pm\SI{0.42}{}$ down to
$\approx\SI{0.056}{\metre\per\second\squared}$ at the same forcing. A precise
free-decay measurement (60 zero-crossings) showed the natural frequency was
\emph{unchanged} at \SI{0.23411}{\hertz}; the effect is therefore \emph{not} a
frequency shift but an increase in the effective damping / a change in the mode
shape (weaker modal projection of the lateral aerodynamic force, and
energy-sharing between the two tower modes) that dropped the forced response by
roughly a factor of seven.

\subsection{Chirp sweep: is there a second SS mode?}
\label{subsec:of-chirp}

To search for a possible second side-to-side tower mode, a linear
wind-direction \emph{chirp} from \SI{0.15}{\hertz} to \SI{2.5}{\hertz} was
applied. The analysis confirmed the first SS mode at \SI{0.234}{\hertz} but
revealed \emph{no} distinct second SS tower mode. The higher-frequency spectral
content was identified as rotor harmonics of the blade-passing frequency: with
the rotor at $\approx\SI{5.7}{rpm}$ the $1P$ frequency is \SI{0.095}{\hertz},
and the observed peaks at \SI{0.283}{}, \SI{0.565}{}, \SI{0.847}{} and
\SI{1.128}{\hertz} correspond to $3P$, $6P$, $9P$ and $12P$. The evenly spaced
side peaks are $1P$ sidebands of the swept forcing against the rotor rotation,
not tower modes. The slow ring-down of the first mode
($\zeta\approx\SI{0.5}{\percent}$, $\tau\approx\SI{136}{\second}$) dominates the
chirp envelope and masks any weaker high-frequency resonance.

\subsection{Discussion and motivation for the reduced model}
\label{subsec:of-discussion}

The OpenFAST experiments establish two firm results. First, the side-to-side
tower mode \emph{can} be driven to a clear, growing resonance, but only through
the \emph{aerodynamic} path (an oscillating wind direction at
$f_{\mathrm{SS}}$); the mechanical response builds up as expected for a lightly
damped mode. Second, the intended \emph{electrical} excitation path is
effectively \emph{closed} in this model: ROSCO computes the generator torque
internally from the measured speed and rejects external power/torque commands,
and the $\texttt{VSContrl}=4$ route that would open the path crashes the FMU.
Consequently the high-fidelity model cannot demonstrate an electrically driven
mechanical resonance without rebuilding the controller interface.

This negative result is the direct motivation for the reduced two-mass model
studied in \S\ref{sec:frequency-analysis}: there the electromagnetic torque
enters the generator-mass equation of motion explicitly through
$T_e = P_e/(\omega_{e,\mathrm{filt}}\eta)$, so the electrical$\to$mechanical
path is \emph{open by construction}. The reduced model therefore provides the
setting in which the electrically driven torsional resonance can actually be
demonstrated (\S\ref{sec:forcedresponse}). Note that the target mode differs
between the two models: the high-fidelity model exhibits a lateral \emph{tower}
mode at \SI{0.234}{\hertz}, whereas the reduced model has no tower and instead
exhibits a \emph{drivetrain torsional} mode at \SI{3.49}{\hertz}. The physical
phenomenon under study --- an electrically driven, lightly damped mechanical
resonance --- is nevertheless the same, and it is the phenomenon, rather than
the specific numerical frequency, that the reduced-model demonstration is
intended to capture.
